\centerline{\Huge Условия и решения задач.}
\centerline{\bf \Huge 6 класс}
 
\problem{6.1}Существует ли дробь, равная $\dfrac{19}{20}$, в которой знаменатель больше числителя ровно на 2023?

\hspace{3mm}

Ответ: $\dfrac{38437}{40460}$.\\

Решение:\\
Заметим: $20\cdot2023 - 19\cdot2023 = 2023\cdot(20-19) = 2023\cdot1 = 2023 $
Домножим числитель и знаменатель дроби на число 2023. Получим $\dfrac{19\cdot2023}{20\cdot2023} = \dfrac{38437}{40460}.$

\hspace{3mm}

Критерии:\\
Приведен пример правильной дроби (дополнительных объяснений не требуется): \textbf{7 баллов.}\\
Идеи, продвигающие ход решения, но недостаточные для получения верного ответа: \textbf{2 балла.}\\
Неправильная дробь или доказательство несуществования: \textbf{0 баллов.}

\hspace{3mm}

\problem{6.2}В каждой из двадцати школ города Элисты обучаются поровну учеников. Если в каждой школе будет обучаться на 50 учеников меньше, то во всех школах останется столько учеников, сколько их раньше было суммарно в восемнадцати школах. Сколько учеников изначально обучалось в каждой школе? 

\hspace{3mm}

Ответ: 500 учеников.

Решение:\\
\textit{I вариант:} Всего во всех школах стало обучаться на 20 · 50 = 1000 учеников меньше. До этого было 20 школ с одинаковым количеством учеников, а стало 18 школ с одинаковым количеством учеников. Значит, общее количество учеников уменьшилось на две школы и в то же время на 1000 учеников.
Следовательно, в каждой школе обучается по 1000/2 = 500 учеников.

\textit{II вариант:} Пусть в каждой школе обучается по X учеников. Тогда составим уравнение:

\newpage

\hspace{3mm}

\noindent
$20\cdot X - 20\cdot50 = 18\cdot X\\
2\cdot X = 20\cdot50 = 1000$\\

\noindent$X = \dfrac{1000}{2} = 500$\\
Следовательно, в каждой школе обучается по 500 учеников.

\hspace{3mm}

Критерии:\\
Любое верное решение не подбором и с верным обоснованием: \textbf{7 баллов.}\\
В целом верное решение с верным обоснованием, но с допущением арифметических ошибок: \textbf{5 баллов.}\\
Решение подбором с обоснованием подбора: \textbf{4 балла.}\\
Идеи, продвигающие ход решения, но недостаточные для получения верного ответа: \textbf{2 балла.}\\
Неверное решение или идеи, не продвигающие ход решения: \textbf{0 баллов.} \\
Верный или неверный ответ без обоснования: \textbf{0 баллов.}

\hspace{3mm}


\problem{6.3}Можно ли покрасить некоторые белые поля доски $2022\times2023$ (изначально все поля белые) в чёрный цвет так, чтобы белых полей было на 2023 меньше, чем чёрных?

\hspace{3mm}

Ответ: Нет, нельзя.

Решение:\\
Предположим, что можно. Тогда пусть белых полей X, а чёрных полей $X + 2023$. Следовательно, суммарное количество полей $X + X + 2023 = 2\cdot X + 2023$, 
где $2\cdot X$ - число чётное при любом натуральном X. 
Значит, число $2\cdot X + 2023$ - нечётное, так как является суммой чётного и нечётного. Получается, общее количество полей на доске нечётно. Но всего полей на доске $2022\times2023$ - чётное количество, так как является произведением чётного и нечётного числа. Противоречие. Следовательно, предположение о том, что это можно было сделать, является неверным. Отсюда следует, что нельзя так покрасить. 

\hspace{3mm}

Критерии:\\
Любое правильное решение с верным обоснованием: \textbf{7 баллов.}\\

\newpage

\hspace{3mm}


\noindentПоказано, что после покраски чёрных и белых суммарно нечетное количество: \textbf{4 балла.}\\
Показано, что всего полей на доске чётное количество: \textbf{2 балла.}\\
Идеи, продвигающие ход решения, но недостаточные для получения верного ответа: \textbf{1 балл.}\\
Неверное решение или идеи, не продвигающие ход решения: \textbf{0 баллов.}\\
Верный или неверный ответ без обоснования: \textbf{0 баллов.}

\hspace{3mm}

\problem{6.4}Три друга - Алдар, Бадма и Хонгор закончили решать городской этап олимпиады по математике, и каждый воскликнул. Алдар: «Я точно умнее Бадмы!», Бадма: «Алдар определенно не самый умный из нас!», Хонгор: «Эх, Алдар умнее меня…». Известно, что только самый умный солгал, а остальные сказали правду. Кто самый умный из друзей, если известно, что такой среди них точно есть?

\hspace{3mm}

Ответ: Хонгор.

Решение:\\
Предположим, что Алдар самый умный. Значит, по условию задачи он солгал, а остальные сказали правду. Но Бадма сказал, что Алдар не самый умный из нас, а он самый умный по предположению. Следовательно, Бадма лжёт, однако лгать может только Алдар. Противоречие. Значит, наше предположение о том, что Алдар самый умный, было неверным. Отсюда следует, что Алдар не самый умный. 

Предположим, что самый умный - Бадма. Значит, по условию задачи он солгал, а остальные сказали правду. Но Алдар сказал, что он умнее Бадмы, а по предположению Бадма самый умный. Значит, Алдар лжёт, но он лгать не может, так как лжёт только один Бадма. Противоречие. Следовательно, наше предположение о том, что Бадма самый умный, было неверным. Таким образом, и Бадма не самый умный.

Мы имеем следующее: Алдар и Бадма не самые умные, а самый умный среди друзей точно есть. Следовательно, самым умным является Хонгор.

\hspace{3mm}

Критерии: \\
Полностью верное решение со всеми обоснованиями: \textbf{7 баллов.}\\
Доказано, что только один из мальчиков (Алдар, Бадма) не самый умный: \textbf{4 балла.}\\

\newpage

\hspace{3mm}

\noindentИдеи, продвигающие ход решения, но недостаточные для получения верного ответа: \textbf{2 балла.}\\
Неверное решение или идеи, не продвигающие ход решения: \textbf{0 баллов.}\\
Верный или неверный ответ без обоснования: \textbf{0 баллов.}

\hspace{3mm}

\problem{6.5}Чингиз записал на доске десятизначное число, которое делится на 6, но его лучший друг Анджа решил над ним подшутить и стёр первую и последнюю цифры этого числа. В результате на доске осталось записано число 20232023. Определите все возможные варианты десятизначного числа, которое мог записать Чингиз до шутки Анджи?

Ответ:\\
1) 1202320230\\
2) 4202320230\\
3) 7202320230\\
4) 2202320232\\
5) 5202320232\\
6) 8202320232\\ 
7) 3202320234\\
8) 6202320234\\
9) 9202320234\\ 
10) 1202320236\\
11) 4202320236\\
12) 7202320236\\ 
13) 2202320238\\
14) 5202320238\\
15) 8202320238

\hspace{3mm}

Решение:\\
Так как число Чингиза делилось на 6, то оно должно делиться на 2 и на 3 одновременно. Следовательно, должны выполняться оба признака делимости на 2 и на 3, а именно: последняя цифра числа должна быть чётной, и сумма цифр числа должна делиться на 3.

\newpage

\hspace{3mm}

\noindentПереберем все варианты:\\
Пусть последняя цифра была 0, тогда сумма всех цифр равна 14. Ближайшие числа, кратные 3, не меньше 14 - это 15, 18, 21, 24 …\\
Для 15 первая цифра должна быть 1.\\
Для 18 первая цифра должна быть 4.\\
Для 21 первая цифра должны быть 7.\\
Для 24 первая цифра должна быть 10, чего не может быть и для вариантов больше также не может быть.\\
Следовательно, могли быть числа:\\
1202320230\\
4202320230\\
7202320230

Пусть последняя цифра была 2, тогда сумма всех цифр равна 16. Ближайшие числа, кратные 3, не меньше 16 - это 18, 21, 24, 27…\\
Для 18 первая цифра должна быть 2.\\
Для 21 первая цифра должна быть 5.\\
Для 24 первая цифра должны быть 8.\\
Для 27 первая цифра должна быть 11, чего не может быть и для вариантов больше также не может быть.\\
Следовательно, могли быть числа:\\
2202320232\\
5202320232\\
8202320232 

Пусть последняя цифра была 4, тогда сумма всех цифр равна 18. Ближайшие числа, кратные 3, не меньше 18 - это 18, 21, 24, 27, 30...\\
Для 18 первая цифра должна быть 0, чего не может быть, так как число не может начинаться с 0.\\
Для 21 первая цифра должна быть 3.\\
Для 24 первая цифра должны быть 6.\\
Для 27 первая цифра должна быть 9.\\
Для 30 первая цифра должна быть 12, чего не может быть и для вариантов больше также не может быть.\\
Следовательно, могли быть числа:

\newpage

\hspace{3mm}

\noindent3202320234\\
6202320234\\
9202320234\\
Пусть последняя цифра была 6, тогда сумма всех цифр равна 20. Ближайшие числа, кратные 3, не меньше 20 - это 21, 24, 27, 30...\\
Для 21 первая цифра должна быть 1.\\
Для 24 первая цифра должны быть 4.\\
Для 27 первая цифра должна быть 7.\\
Для 30 первая цифра должна быть 10, чего не может быть и для вариантов больше также не может быть.\\
Следовательно, могли быть числа:\\
1202320236\\
4202320236\\
7202320236 

Пусть последняя цифра была 8, тогда сумма всех цифр равна 22. Ближайшие числа, кратные 3, не меньше 22 - это 24, 27, 30, 33...\\
Для 24 первая цифра должна быть 2.\\
Для 27 первая цифра должны быть 5.\\
Для 30 первая цифра должна быть 8.\\
Для 33 первая цифра должна быть 11, чего не может быть и для вариантов больше также не может быть.\\
Следовательно, могли быть числа:\\
2202320238\\
5202320238\\
8202320238 

Выпишем все возможные варианты:\\
1) 1202320230\\
2) 4202320230\\
3) 7202320230\\
4) 2202320232\\
5) 5202320232\\
6) 8202320232\\ 
7) 3202320234

\newpage

\hspace{3mm}


\noindent8) 6202320234\\
9) 9202320234\\ 
10) 1202320236\\
11) 4202320236\\
12) 7202320236\\ 
13) 2202320238\\
14) 5202320238\\
15) 8202320238 

\hspace{3mm}

Критерии:\\
15 правильных чисел со всеми обоснованиями: \textbf{7 баллов.}\\
10-14 правильных чисел со всеми обоснованиями: \textbf{6 баллов.}\\
15 правильных чисел без обоснований: \textbf{5 баллов.}\\
10-14 правильных чисел без обоснований: \textbf{4 балла.}\\
5-9 правильных чисел со всеми обоснованиями: \textbf{4 балла.}\\
5-9 правильных чисел без обоснований: \textbf{3 балла.}\\
1-4 правильных чисел с обоснованием: \textbf{2 балла.}\\
1-4 правильных чисел без обоснования: \textbf{1 балл.}\\
0 правильных чисел с хорошими идеями, помогающие решению: \textbf{1 балл.}\\
0 правильных чисел с идеями, не помогающие решению или неправильное решение: \textbf{0 баллов.}
